% gen_1
\begin{lstlisting}[language=Diff]
diff --git a/task_agent.py b/task_agent.py
index 3798256..98e6706 100644
--- a/task_agent.py
+++ b/task_agent.py
@@ -5,7 +5,7 @@ from utils.common import extract_jsons
 class TaskAgent(AgentSystem):
     def forward(self, inputs):
         """
-        An agent that solves a given task.
+        An agent that solves a given task with enhanced reasoning and error handling.
         ...
         """
-        domain = inputs['domain']
-        instruction = f"""You are an agent.
+        domain = inputs.get('domain', 'general')
+
+        # Enhanced instruction with clearer structure and reasoning guidance
+        instruction = f"""You are an expert AI agent specialized in solving complex tasks.

-Task input:
+Task Domain: {domain}
+
+Task Input:
 ```
 {inputs}
 ```

+Instructions:
+1. Carefully analyze the task input and identify the key requirements
+2. Break down the problem into logical steps if needed
+3. Formulate your response based on the task requirements
+4. Provide your final answer in the specified JSON format
+
 Respond in JSON format with the following schema:
 <json>
 {{
-    "response": ...
+    "reasoning": "Brief explanation of your approach",
+    "response": "Your final answer here"
 }}
-</json>"""
-        new_msg_history = chat_with_agent(instruction, model=self.model, msg_history=[], logging=self.log)
+</json>

-        # Extract the response
+IMPORTANT: Always include both 'reasoning' and 'response' fields in your JSON output."""
+
+        # Attempt to get response with retry logic
+        max_retries = 2
         prediction = "None"
-        try:
-            extracted_jsons = extract_jsons(new_msg_history[-1]['text'])
-            if extracted_jsons is not None and "response" in extracted_jsons[-1]:
-                prediction = extracted_jsons[-1]['response']
-        except Exception as e:
-            self.log(f"Error extracting prediction: {e}")
-            prediction = "None"
+        new_msg_history = []
+
+        for attempt in range(max_retries + 1):
+            try:
+                new_msg_history = chat_with_agent(...)
+                # Extract the response with retry logic
+                if new_msg_history and len(new_msg_history) > 0:
+                    extracted_jsons = extract_jsons(new_msg_history[-1]['text'])
+                    if extracted_jsons is not None and len(extracted_jsons) > 0:
+                        response_json = extracted_jsons[-1]
+                        if "reasoning" in response_json:
+                            self.log(f"Agent reasoning: {response_json['reasoning']}")
+                        if "response" in response_json:
+                            prediction = response_json['response']
+                            break
+            except Exception as e:
+                self.log(f"Error in attempt {attempt + 1}: {str(e)}")

         return prediction, new_msg_history
\end{lstlisting}

% % gen_7
% \begin{lstlisting}[language=Diff]
% diff --git a/task_agent.py b/task_agent.py
% index 98e6706..3c2dd8a 100644
% --- a/task_agent.py
% +++ b/task_agent.py
% @@ -17,8 +17,43 @@ class TaskAgent(AgentSystem):
%          """
%          domain = inputs.get('domain', 'general')

% -        # Enhanced instruction with clearer structure and reasoning guidance
% -        instruction = f"""You are an expert AI agent specialized in solving complex tasks.
% +        # Build domain-specific instruction
% +        if domain == 'paper_review':
% +            paper_text = inputs.get('paper_text', '')
% +            instruction = f"""You are an expert reviewer for a top-tier academic conference.

% +Your task is to decide whether to ACCEPT or REJECT the following paper.

% +Paper Text:
% +{paper_text}

% +Evaluation Criteria (consider ALL criteria):
% +1. **Novelty**: Does the paper present genuinely new ideas or significant contributions?
% +2. **Technical Quality**: Is the methodology sound? Are experiments well-designed?
% +3. **Clarity**: Is the paper well-written and easy to understand?
% +4. **Significance**: Will this work meaningfully impact the field?
% +5. **Validation**: Are claims properly supported by evidence?

% +Instructions:
% +1. Analyze the paper objectively against all five criteria
% +2. Identify both strengths and weaknesses
% +3. Make a balanced decision - avoid default bias toward acceptance or rejection
% +4. Base your decision on concrete evidence from the paper
% +5. Accept papers that meet high standards across most criteria
% +6. Reject papers with fatal flaws or insufficient contributions

% +Respond in JSON format:
% +<json>
% +{{
% +    "reasoning": "Balanced analysis of key strengths and weaknesses (2-4 sentences)",
% +    "response": "accept or reject (exactly one word, lowercase)"
% +}}
% +</json>

% +IMPORTANT: The response field must contain ONLY "accept" or "reject" (lowercase)."""
% +        else:
% +            # Generic instruction for other domains
% +            instruction = f"""You are an expert AI agent specialized in solving complex tasks.
% \end{lstlisting}

% gen_8
\begin{lstlisting}[language=Diff]
diff --git a/task_agent.py b/task_agent.py
index 3c2dd8a..faab41d 100644
--- a/task_agent.py
+++ b/task_agent.py
@@ -27,30 +27,96 @@ Paper Text:
+Critical Evaluation Guidelines:
+- Be rigorous and selective - most submissions have issues
+- Accept ONLY papers that are strong across MOST criteria (4-5/5)
+- Reject papers with significant weaknesses in multiple areas
+- A paper needs to be clearly above average to warrant acceptance
+- Consider: Would this paper strengthen the conference program?
+- If a paper is borderline or has major concerns, lean toward REJECT

+        elif domain == 'genesis_go2walking':
+            # Genesis domain-specific instruction with environment attribute documentation
+            task_description = inputs.get('task_description', '')
+            genesis_env_path = inputs.get('genesis_environment_path', '')
+            default_reward_function = inputs.get('default_reward_function', '')
+
+            instruction = f"""You are an expert in reinforcement learning and reward function design.

+Task: {task_description}

+Default Reward Function:
+```python
+{default_reward_function}
+```

+AVAILABLE ENVIRONMENT ATTRIBUTES (use ONLY these):
+- env.commands: [num_envs, 3] - velocity commands [lin_vel_x, lin_vel_y, ang_vel_z]
+- env.base_lin_vel: [num_envs, 3] - base linear velocity [x, y, z]
+- env.base_ang_vel: [num_envs, 3] - base angular velocity [roll, pitch, yaw]
+- env.base_pos: [num_envs, 3] - base position [x, y, z]
+- env.base_quat: [num_envs, 4] - base orientation quaternion
+- env.projected_gravity: [num_envs, 3] - gravity vector in base frame
+- env.dof_pos: [num_envs, 12] - joint positions
+- env.dof_vel: [num_envs, 12] - joint velocities
+- env.actions: [num_envs, 12] - current actions
+- env.last_actions: [num_envs, 12] - previous actions
+- env.last_dof_vel: [num_envs, 12] - previous joint velocities

+IMPORTANT CONSTRAINTS:
+- DO NOT use env.torques (not available)
+- DO NOT access undefined attributes
+- Use only the attributes listed above
+- All operations must work with PyTorch tensors
\end{lstlisting}

% gen_12
\begin{lstlisting}[language=Diff]
diff --git a/task_agent.py b/task_agent.py
index faab41d..62c444f 100644
--- a/task_agent.py
+++ b/task_agent.py
@@ -88,19 +90,24 @@ AVAILABLE ENVIRONMENT ATTRIBUTES (use ONLY these):
-IMPORTANT CONSTRAINTS:
-- DO NOT use env.torques (not available)
-- DO NOT access undefined attributes
-- Use only the attributes listed above
-- All operations must work with PyTorch tensors
+CRITICAL RULES (VIOLATION WILL CAUSE IMMEDIATE FAILURE):
+1. DO NOT reference env.torques or env.dt in your code - THESE DO NOT EXIST
+2. DO NOT add hasattr() checks for unavailable attributes and then use them
+3. DO NOT assume any attributes beyond the exact list above
+4. NEVER write code like: if hasattr(env, 'torques'): use env.torques
+5. All operations must work with PyTorch tensors
+6. If you need time differences, use fixed timestep: dt = 0.02
+7. ONLY use attributes from the AVAILABLE list - nothing else exists

+                            # Post-process prediction based on domain
+                            if domain == 'paper_review':
+                                # Ensure response is exactly "accept" or "reject"
+                                prediction_lower = str(prediction).lower().strip()
+                                if 'accept' in prediction_lower and 'reject' not in prediction_lower:
+                                    prediction = 'accept'
+                                elif 'reject' in prediction_lower:
+                                    prediction = 'reject'
+                                else:
+                                    self.log(f"Warning: Ambiguous response '{prediction}', defaulting to 'reject'")
+                                    prediction = 'reject'
\end{lstlisting}

% gen_13
\begin{lstlisting}[language=Diff]
diff --git a/task_agent.py b/task_agent.py
index 62c444f..fb504a2 100644
--- a/task_agent.py
+++ b/task_agent.py
@@ -40,26 +40,30 @@ Critical Evaluation Guidelines:
-- If a paper is borderline or has major concerns, lean toward REJECT
+- Balance your decisions: aim for appropriate selectivity (not too lenient, not too harsh)
+- If a paper is borderline, carefully weigh the evidence on both sides

 Reward Function Design Guidelines:
 1. Primary reward: Track forward velocity command (env.commands[:, 0] vs env.base_lin_vel[:, 0])
+   - Use exponential reward: torch.exp(-error * temperature) for smooth gradients
+   - Typical temperature: 1.0-2.0, scale: 0.8-1.0
 2. Secondary rewards: Track angular velocity, maintain base height, stable orientation
-3. Penalties: Excessive action changes (use env.actions - env.last_actions)
-4. Use exponential rewards: torch.exp(-error * temperature) for smooth gradients
-5. Scale rewards appropriately (main objectives: 0.5-1.0, penalties: -0.5 to -0.01)
-6. For smooth action changes: use torch.sum(torch.square(env.actions - env.last_actions))
-7. For velocity smoothness: use torch.sum(torch.square(env.dof_vel - env.last_dof_vel))
-8. Return: total_reward, reward_components dict, reward_scales dict
+   - Angular velocity tracking: similar to linear velocity
+   - Base height: penalize deviation from target (typically 0.32m)
+   - Orientation: use env.projected_gravity to penalize tilt
+3. Penalties for smooth and stable locomotion:
+   - Action smoothness: -0.01 * torch.sum(torch.square(env.actions - env.last_actions), dim=1)
+   - Joint velocity smoothness: -0.001 * torch.sum(torch.square(env.dof_vel - env.last_dof_vel), dim=1)
+   - Excessive joint velocities: -0.0005 * torch.sum(torch.square(env.dof_vel), dim=1)
+   - Unwanted lateral movement: penalize env.base_lin_vel[:, 1]
+4. Balance reward components:
+   - Main tracking objectives: 0.5-1.0 scale
+   - Secondary objectives: 0.2-0.5 scale
+   - Smoothness penalties: -0.01 to -0.001 scale
+5. Return: total_reward, reward_components dict, reward_scales dict
\end{lstlisting}

% % gen_14
% \begin{lstlisting}[language=Diff]
% diff --git a/task_agent.py b/task_agent.py
% index fb504a2..3de640e 100644
% --- a/task_agent.py
% +++ b/task_agent.py
% @@ -35,35 +35,38 @@ Evaluation Criteria (consider ALL criteria carefully):

%  Critical Evaluation Guidelines:
% -- Be rigorous and selective - most submissions have issues
% -- Accept ONLY papers that are strong across MOST criteria (4-5/5)
% -- Reject papers with significant weaknesses in multiple areas
% +- Be rigorous and selective - top conferences accept ~25% of submissions
% +- Most papers have significant issues that warrant rejection
% +- Accept ONLY papers that excel in novelty AND technical quality AND significance
% +- Even well-written papers should be rejected if they lack novelty or significance
% +- Papers must have STRONG validation with comprehensive experiments
% +- If a paper has ANY critical flaws in methodology or lacks clear novelty, REJECT
% +- For borderline cases: When in doubt between accept/reject, favor REJECT
% +- Ask yourself: Is this among the top 25% of papers? If not clearly yes, then REJECT

%  Instructions:
% -1. Critically analyze the paper against all five criteria
% -2. Identify both strengths and weaknesses objectively
% -3. Be honest about limitations - most papers have them
% +1. Start by identifying potential RED FLAGS: lack of novelty, weak validation, methodological flaws
% +2. Critically analyze against ALL five criteria - weakness in ANY major criterion -> likely REJECT
% +3. For each criterion, explicitly note whether it meets the HIGH bar for a top conference
% +4. List specific weaknesses clearly - be thorough in finding problems
% +5. Only if the paper is STRONG on novelty AND technical quality AND significance should you consider accept
% +6. Apply the "top 25% test": Would this paper be in the top quartile of all submissions?
% +7. Default stance should be critical - acceptance requires exceptional quality

%  IMPORTANT:
%  - The response field must contain ONLY "accept" or "reject" (lowercase, no other text)
% -- Provide detailed reasoning that justifies your decision
% -- Consider all five evaluation criteria in your analysis"""
% +- Start reasoning by identifying weaknesses - be thorough and critical
% +- Accept ONLY if paper excels across ALL major criteria (novelty, technical quality, significance)
% +- When uncertain or borderline, default to REJECT
% +- Remember: Most papers should be rejected; acceptance is for exceptional work only"""
% \end{lstlisting}

% % gen_16
% \begin{lstlisting}[language=Diff]
% diff --git a/task_agent.py b/task_agent.py
% index 3de640e..b39af37 100644
% --- a/task_agent.py
% +++ b/task_agent.py
% @@ -20,53 +20,49 @@ class TaskAgent(AgentSystem):
%          if domain == 'paper_review':
%              paper_text = inputs.get('paper_text', '')
% -            instruction = f"""You are an expert reviewer for a top-tier academic conference.
% +            instruction = f"""You are a reviewer for a top-tier academic conference with a ~50% acceptance rate.

% +**Evaluation Process:**

% +Score each criterion on a 3-point scale:
% +- **3 = Strong**: Exceeds expectations significantly
% +- **2 = Adequate**: Meets minimum standards but nothing exceptional
% +- **1 = Weak**: Below acceptable standards

% +**Five Criteria:**
% +1. **Novelty**: Does it present genuinely new ideas?
% +2. **Technical Quality**: Is methodology sound and rigorous?
% +3. **Clarity**: Is it well-written and clear?
% +4. **Significance**: Will it meaningfully impact the field?
% +5. **Validation**: Are claims properly validated with evidence?

% +**Decision Rules (FOLLOW STRICTLY):**
% +- **ACCEPT**: Average score >= 2.4 AND no scores of 1 on Novelty, Technical Quality, or Significance
% +- **REJECT**: Average score < 2.4 OR any score of 1 on Novelty, Technical Quality, or Significance

% +**Examples:**
% +- Scores [3,3,2,3,2] -> Average 2.6 -> ACCEPT
% +- Scores [3,2,3,2,2] -> Average 2.4 -> ACCEPT
% +- Scores [2,2,3,2,2] -> Average 2.2 -> REJECT
% +- Scores [3,1,3,3,3] -> Has 1 on Technical Quality -> REJECT
% +- Scores [2,3,2,2,2] -> Average 2.2 -> REJECT
% \end{lstlisting}

% % gen_18
% \begin{lstlisting}[language=Diff]
% diff --git a/task_agent.py b/task_agent.py
% index b39af37..7a30be1 100644
% --- a/task_agent.py
% +++ b/task_agent.py
% @@ -20,49 +20,67 @@ class TaskAgent(AgentSystem):
%          if domain == 'paper_review':
%              paper_text = inputs.get('paper_text', '')
% -            instruction = f"""You are a reviewer for a top-tier academic conference with a ~50% acceptance rate.
% +            instruction = f"""You are a CRITICAL reviewer for a top-tier conference (50% acceptance rate). Be skeptical and demand excellence.

% +**CRITICAL SCORING RUBRIC** (Be strict - most papers are NOT excellent):

% +**Score 3 (Strong)** - RARE, only for exceptional work:
% +- Novelty: Introduces fundamentally new approach/paradigm that will change how people think
% +- Technical: Methodology is impeccable with deep theoretical insights or rigorous proofs
% +- Clarity: Exceptionally clear writing that could be a textbook example
% +- Significance: Will spawn new research direction or solve major open problem
% +- Validation: Comprehensive experiments covering all edge cases with statistical rigor

% +**Score 2 (Adequate)** - Barely acceptable for publication:
% +- Novelty: Minor incremental improvement or simple combination of existing ideas
% +- Technical: Methodology is correct but straightforward, no major technical contributions
% +- Clarity: Readable but could be clearer, some concepts need better explanation
% +- Significance: Addresses niche problem with limited broader impact
% +- Validation: Basic experiments shown, but missing important ablations or baselines

% +**Score 1 (Weak)** - Below publication threshold:
% +- Novelty: No new ideas, just applies existing methods to new dataset
% +- Technical: Flawed methodology, missing baselines, or questionable assumptions
% +- Clarity: Hard to understand, poorly organized, or missing critical details
% +- Significance: Problem is trivial or solution won't impact the field
% +- Validation: Insufficient experiments, cherry-picked results, or weak evaluation

% +**RED FLAGS (Consider score 1 if ANY present):**
% +- Claims not supported by experiments
% +- Missing comparisons to obvious baselines
% +- Unclear methodology that cannot be reproduced
% +- Overclaimed significance without evidence
% +- Fundamental flaws in approach
% +- Incremental work (just benchmarking existing methods)

% +**DECISION RULES (FOLLOW EXACTLY):**
% +- **ACCEPT**: Average >= 2.5 AND all critical criteria (Novelty, Technical, Significance) >= 2
% +- **REJECT**: Average < 2.5 OR ANY of (Novelty, Technical, Significance) = 1

% +**REMEMBER:**
% +- Be CRITICAL - most papers should get 2s, not 3s
% +- Score 2 means "barely adequate", not "good"
% +- When in doubt between scores, choose the lower one
% +- ~50% of papers should be rejected"""
% \end{lstlisting}

% % gen_37
% \begin{lstlisting}[language=Diff]
% diff --git a/task_agent.py b/task_agent.py
% index 7a30be1..42c0f36 100644
% --- a/task_agent.py
% +++ b/task_agent.py
% @@ -25,48 +25,50 @@ Paper Text:

% -**CRITICAL SCORING RUBRIC** (Be strict - most papers are NOT excellent):
% +**BALANCED SCORING RUBRIC** (Be fair and objective):

% -**Score 3 (Strong)** - RARE, only for exceptional work:
% +**Score 3 (Strong)** - Clearly exceeds standards:
% +- Novelty: Presents genuinely new ideas, significant advancement over prior work
% +- Technical: Sound methodology with rigorous validation, thorough analysis
% +- Clarity: Well-written, clear presentation, reproducible methodology
% +- Significance: Addresses important problem, results will impact the field
% +- Validation: Comprehensive experiments with proper baselines and analysis

% -**Score 2 (Adequate)** - Barely acceptable for publication:
% +**Score 2 (Adequate)** - Meets publication standards:
% +- Novelty: Incremental but meaningful contribution, builds reasonably on prior work
% +- Technical: Methodology is sound though not groundbreaking, adequate rigor
% +- Clarity: Generally clear though some parts could be better explained
% +- Significance: Addresses relevant problem with reasonable impact
% +- Validation: Sufficient experiments to support claims, covers main scenarios

% -**Score 1 (Weak)** - Below publication threshold:
% +**Score 1 (Weak)** - Below publication standards:

% +**CRITICAL REJECTION CRITERIA:**
% +Reject if ANY of these are present:
% +- Fatal technical flaws or incorrect methodology
% +- Claims clearly not supported by evidence
% +- Cannot be reproduced due to missing details
% +- Novelty score = 1 (no new contribution)
% +- Technical score = 1 (fundamentally flawed)
% +- Significance score = 1 (trivial problem/solution)

% -**DECISION RULES (FOLLOW EXACTLY):**
% -- **ACCEPT**: Average >= 2.5 AND all critical criteria >= 2
% -- **REJECT**: Average < 2.5 OR ANY critical = 1
% +**DECISION RULES (FOLLOW STRICTLY):**
% +- **ACCEPT**: Average >= 2.4 AND no score of 1 on Novelty, Technical, or Significance
% +- **REJECT**: Average < 2.4 OR ANY score of 1 on Novelty, Technical, or Significance

% +**IMPORTANT REMINDERS:**
% +- Be objective and fair - evaluate each criterion independently
% +- Score 2 means "adequate/acceptable", score 3 means "strong/excellent"
% +- Focus on identifying fatal flaws (score 1) that warrant rejection
% +- Consider both strengths and weaknesses before deciding
% +- Aim for ~50% acceptance rate by being discerning but not overly harsh"""
% \end{lstlisting}

% % gen_43
% \begin{lstlisting}[language=Diff]
% diff --git a/task_agent.py b/task_agent.py
% index 42c0f36..c1c37a5 100644
% --- a/task_agent.py
% +++ b/task_agent.py
% @@ -20,70 +20,88 @@ class TaskAgent(AgentSystem):
%          if domain == 'paper_review':
%              paper_text = inputs.get('paper_text', '')
% -            instruction = f"""You are a CRITICAL reviewer for a top-tier conference (50% acceptance rate).
% +            instruction = f"""You are a RIGOROUS reviewer for a top-tier ML conference with a 25-30% acceptance rate. Most submissions have merit but only the best should be accepted.

% +**REJECTION CRITERIA** (Reject if ANY apply):
% +1. Average score < 2.5 (below quality threshold)
% +2. ANY dimension scores 1 (critical weakness)
% +3. Novelty <= 1.5 (insufficient contribution)
% +4. Technical quality = 1 (fundamentally flawed)
% +5. Claims not supported by evidence
% +6. Cannot be reproduced or validated

% +**ACCEPTANCE CRITERIA** (ALL must be satisfied):
% +1. Average score >= 2.5
% +2. NO score of 1 in any dimension
% +3. Novelty >= 2 (meaningful contribution)
% +4. Technical >= 2 (sound methodology)
% +5. Claims properly supported

% +**REVIEW PROCESS:**
% +1. Carefully read the paper and identify its main contribution
% +2. Score EACH dimension independently (don't let one dimension bias others)
% +3. Be CRITICAL: Look for weaknesses, gaps, missing comparisons
% +4. Check if rejection criteria are met (if yes, must reject)
% +5. Check if ALL acceptance criteria are met (if not all, must reject)
% +6. Calculate average and make final decision

% +**IMPORTANT:**
% +- Default to REJECT unless paper clearly meets ALL acceptance criteria
% +- ~70% of papers should be rejected (maintain high standards)
% +- Score 2 = "acceptable but not great", Score 3 = "genuinely strong"
% +- Most papers will have mix of 2s and 3s, with occasional 1s
% +- Don't be overly generous - top conferences reject many good papers
% \end{lstlisting}

% % gen_89
% \begin{lstlisting}[language=Diff]
% diff --git a/task_agent.py b/task_agent.py
% index c1c37a5..bacc897 100644
% --- a/task_agent.py
% +++ b/task_agent.py
% @@ -20,88 +20,73 @@ class TaskAgent(AgentSystem):
%          if domain == 'paper_review':
%              paper_text = inputs.get('paper_text', '')
% -            instruction = f"""You are a RIGOROUS reviewer for a top-tier ML conference with a 25-30% acceptance rate.
% +            instruction = f"""You are a reviewer for a top-tier ML/AI conference with approximately 50% acceptance rate.

% +**Review Philosophy:**
% +- Be fair and balanced: recognize both strengths and weaknesses
% +- Most papers submitted to top venues have merit - distinguish publishable from not-yet-ready
% +- Use the full 1-3 scale: avoid clustering scores (distribute naturally)
% +- Score 2 is the "default acceptable" level - most criteria will be 2 or 3
% +- Reserve score of 1 for truly serious flaws only

% +**Scoring Scale (Be Generous but Fair):**
% +- **3 = Good/Strong**: Clear strengths on this criterion, above average quality
% +- **2 = Acceptable/Adequate**: Meets publication standards, no major issues
% +- **1 = Poor/Inadequate**: Serious problems that disqualify the paper

% +**Five Evaluation Criteria (score each 1-3):**
% +1. **Novelty**: Is the contribution sufficiently novel? (2=incremental but solid, 3=novel idea)
% +2. **Technical Quality**: Is methodology sound? (2=adequate rigor, 3=excellent execution)
% +3. **Clarity**: Is it well-written? (2=understandable with minor issues, 3=very clear)
% +4. **Significance**: Does it matter? (2=useful contribution, 3=important impact)
% +5. **Validation**: Are claims supported? (2=adequate evidence, 3=thorough validation)

% +**Decision Rules (Apply EXACTLY):**
% +1. Calculate average of all 5 scores
% +2. Check for any fatal flaws (score of 1 on Novelty, Technical Quality, OR Significance)
% +3. **ACCEPT** if: Average >= 2.2 AND no fatal flaws (score of 1 on critical criteria)
% +4. **REJECT** if: Average < 2.2 OR has fatal flaw on critical criterion

% +**Scoring Calibration (Target 50% acceptance):**
% +ACCEPT examples:
% +- [3,3,3,2,2] -> Avg=2.6 -> ACCEPT (strong paper)
% +- [3,2,2,3,2] -> Avg=2.4 -> ACCEPT (good overall)
% +- [2,3,2,2,2] -> Avg=2.2 -> ACCEPT (meets threshold)
% +- [2,2,3,2,2] -> Avg=2.2 -> ACCEPT (adequate)
% +- [2,2,2,3,2] -> Avg=2.2 -> ACCEPT (borderline but acceptable)

% +REJECT examples:
% +- [2,2,2,2,1] -> Avg=1.8 -> REJECT (below threshold)
% +- [2,2,1,2,2] -> Avg=1.8 -> REJECT (below threshold)
% +- [3,1,3,3,2] -> REJECT (fatal flaw: Technical Quality = 1)
% +- [1,3,3,2,3] -> REJECT (fatal flaw: Novelty = 1)
% +- [2,2,2,2,2] -> Avg=2.0 -> REJECT (below threshold)

% +**Critical Guidance:**
% +- **Default to score of 2** for most criteria unless there's a clear reason for 1 or 3
% +- Score 3: Clear strength or excellence on this dimension
% +- Score 2: Acceptable quality, publishable standard (MOST COMMON)
% +- Score 1: Major flaw or serious inadequacy (RARE - use sparingly)
% +- Don't look for perfection - accept solid, publishable work
% +- The 2.2 threshold with proper scoring yields ~50% acceptance
% +- Be generous with 2s and 3s; reserve 1s for truly poor work

% +IMPORTANT:
% +- Most papers should get mostly 2s and 3s (score 1 is rare)
% +- Calculate average correctly (sum/5)
% +- Apply decision rules exactly: avg >= 2.2 AND no fatal flaws -> ACCEPT
% +- Response field: ONLY "accept" or "reject" (lowercase, no other text)"""
% \end{lstlisting}
